\section*{Hoofdstuk \ref{ch:b2:two-wave-interference} --- Tweegolfsinterferentie}

\begin{solution}{exo:b2:two-wave-interference:1}
$i = \lambda D/a = \qty{2.4}{mm}$; $x/i = 4.2$; tiende donkere streep bij
$9.5i = \qty{22.5}{mm}$;
$a = \qty{2}{mm}$: \qty{0.47}{mm}; in water $\lambda/n$: \qty{1.8}{mm}.
\end{solution}

\begin{solution}{exo:b2:two-wave-interference:2}
$1.25I_0 \pm I_0$: $V = 0.8$. $1.01I_0 \pm 0.2I_0$: $V = 0.2$ --- de
interferentieterm
gaat als de verhouding van de \emph{amplituden}, $1/10$, niet die van de
intensiteiten. Filter:
$V = 2\sqrt{0.5}/1.5 = 0.94$.
\end{solution}

\begin{solution}{exo:b2:two-wave-interference:3}
$2ne = \qty{665}{nm}$: helder bij $665/(p + \tfrac12)$: \qty{443}{nm} (blauwviolet);
donker bij $665/p$: \qty{665}{nm} (rood): een blauw vlies.
$e = \qty{1}{\micro m}$: $2ne = \qty{2660}{nm}$,
helder bij $760$, $591$, $484$, \qty{409}{nm} --- vier kleuren tegelijk, een
verwassen rozeachtig wit.
\end{solution}

\begin{solution}{exo:b2:two-wave-interference:4}
$r_p = \sqrt{p\lambda R}$: $1.09$, $1.53$, \qty{1.88}{mm};
$p = r^2/\lambda R = 85$ ringen binnen
\qty{1}{cm}; het midden is donker: dikte nul en één tekenomkering. Water:
stralen gedeeld door $\sqrt n$, $98$ ringen.
\end{solution}

\begin{solution}{exo:b2:two-wave-interference:5}
(a) $i = 0.8$ en \qty{1.4}{mm}; tweede rode \qty{2.8}{mm}, derde violette
\qty{2.4}{mm}.
(b) Orde $p$ rood bij $1.4p$ mm ontmoet orde $p + 1$ violet bij $0.8(p + 1)$: vanaf
$p = 2$ overlappen de orden. (c) $\delta = ax/D = \qty{2.5}{\micro m}$:
ontbrekende $\lambda = \delta/(p +
\tfrac12)$: $714$, $556$, \qty{455}{nm}.
(d) Wit midden; dan strepen met een
blauwe binnenrand en een rode buitenrand; tegen de derde orde mengen de kleuren
zich tot wit.
\end{solution}

\begin{solution}{exo:b2:two-wave-interference:6}
(a) $b < \lambda L/4a = \qty{59}{\micro m}$. (b) $\pi ab/\lambda L = 6.7$:
contrast $|\sin6.7|/6.7
= 0.06$. (c) Een bron van een centimeter breed
beslaat honderden streepafstanden
verschuiving: geen strepen. (d) $a < \lambda/\theta = \qty{61}{\micro m}$.
\end{solution}

\begin{solution}{exo:b2:two-wave-interference:7}
(a) $a = 2h = \qty{0.4}{mm}$: $i = \lambda D/a = \qty{1.25}{mm}$. (b) Bij
$x = 0$ zijn de wegen
gelijk en is de streep toch donker: de weerkaatsing voegt $\pi$ toe (tekenomkering
bij scherende inval op het dichtere midden). (c)
$\delta = ax/D <
\qty{50}{\micro m}$: $100$ strepen. (d) De zee is de
spiegel, de antenne op de
klip het scherm: terwijl de bron rijst, trekken de strepen voorbij en
geven zij haar hoogte.
\end{solution}

\begin{solution}{exo:b2:two-wave-interference:8}
(a) $2d = 42\lambda$: $d = \qty{12.4}{\micro m}$. (b)
$100/42 = \qty{2.4}{mm}$. (c) $2nd = p\lambda$:
$56$ strepen. (d) De strepen bewegen naar het verre uiteinde en spreiden
zich terwijl $d$
daalt; een kwart streep is $\lambda/8 = \qty{74}{nm}$ op $d$.
\end{solution}

\begin{solution}{exo:b2:two-wave-interference:9}
(a) Alleen de eerste (lucht $\to$ olie, dichter); olie $\to$ water gaat naar
een lagere
index: $\delta = 2ne + \lambda/2$, helder bij $2ne = (p + \tfrac12)\lambda$.
(b) $2ne =
\qty{1160}{nm}$: versterkt $773$ (rand van het zichtbare),
\qty{464}{nm} (blauw);
uitgedoofd $1160/p$: \qty{580}{nm} (geel): blauw. (c) $2ne = \lambda/2$:
$e = \qty{72}{nm}$.
(d) $\delta = 2ne\cos r$ daalt met de hoek: de versterkte golflengten
schuiven naar het blauw. (e) Dan keren beide weerkaatsingen het teken om, geen
extra $\lambda/2$: helder bij $2ne = p\lambda$, donker bij
$(p + \tfrac12)\lambda$ --- de
\omterm{prop:b2:wave-interfaces:optics}{antireflectielaag} van een kwartgolf.
\end{solution}

\begin{solution}{exo:b2:two-wave-interference:10}
(a) $(n - 1)e = \qty{10}{\micro m}$, $20$ strepen, naar de bedekte spleet toe (waar
de meetkundige weg \qty{10}{\micro m} korter moet worden om $\delta = 0$ te
herstellen):
$x = (n - 1)eD/a = \qty{2}{cm}$. (b) De witte streep ligt daar; de dispersie
van het glas
($\Delta n \approx 0.01$ over het zichtbare, dat wil zeggen \qty{0.2}{\micro m} weg)
blijft onder $\ell_c$: nog zichtbaar, lichtjes gekleurd. (c) Meet de
verschuiving van de witte streep: $e = (\text{verschuiving})\,a/(n - 1)D$.
(d) $e(n - 1)\theta^2
(1 + 1/n)/2 = \qty{63}{nm}$ bij \qty{5}{\degree}: een
achtste streep.
\end{solution}

\begin{solution}{exo:b2:two-wave-interference:11}
(a) $\theta = \qty{2.3e-7}{rad}$: $a = 1.22\lambda/\theta = \qty{3.1}{m}$.
(b) Zijn grens $1.22\lambda/d =
\qty{0.057}{''}$ ligt boven de diameter van
de ster, en de dampkring vervaagt
tot $1''$. (c) $1.22\lambda/100 = \qty{7e-9}{rad} = \qty{1.4}{mas}$. (d) De
strepen vergen
$\delta < \ell_c = \lambda^2/\Delta\lambda = \qty{32}{\micro m}$.
\end{solution}

\begin{solution}{exo:b2:two-wave-interference:12}
(a) $\delta = 2ne\cos r + \lambda/2$;
$\cos r = \sqrt{1 - \sin^2i/n^2} \approx 1 - i^2/2n^2$.
(b) $2ne/\lambda = 10186.8$: het midden, met $p + \tfrac12$ tussen $10186.5$ en
$10187.5$, is noch helder noch donker. (c)
$i_p^2 = 2n^2[1 - (p + \tfrac12)\lambda/
2ne]$ voor $p = 10186$, $10185$,
$10184$: $i = 0.011$, $0.024$, \qty{0.032}{rad}, $r = fi
= 5.4$, $11.8$,
\qty{15.8}{mm}. (d) De hoek $i$ alleen legt de orde vast: elk
bronpunt voedt dezelfde ring onder dezelfde hoek, zodat de ringen, op
oneindig gezien, een brede bron overleven.
\end{solution}

\begin{solution}{pb:b2:two-wave-interference:1}
\textbf{1.} Elke spiegel beeldt $S$ door symmetrie af op de afstand $d$ van
de raaklijn; de twee beelden liggen op de cirkel met straal $d$, op
de hoek $2\alpha$ uit elkaar: $a = 2\alpha d = \qty{2}{mm}$.

\textbf{2.} $S_1$ en $S_2$ zijn twee beelden van dezelfde trilling: hun
fasesprongen zijn dezelfde.

\textbf{3.}
$i = \lambda(d + D)/a = 589 \times 10^{-9} \times 1.7/2 \times 10^{-3} = \qty{0.50}{mm}$:
tien
strepen over \qty{5}{mm}.

\textbf{4.} $\delta = ax/(d + D) = \qty{2.9}{\micro m}$ aan de rand,
$p = 5$: ver onder
$\ell_c/\lambda = 34\,000$.

\textbf{5.} Verschuiving
$bD'/d = 20 \times 10^{-6} \times 1.7/0.2 = \qty{0.17}{mm}$, een derde van
$i$, zodat het contrast wordt
$|\operatorname{sinc}(\pi ab/\lambda d)| = \operatorname{sinc}(1.07) = 0.82$.

\textbf{6.} $\pi ab/\lambda d = 5.3$: $|\sin5.3|/5.3 = 0.16$.

\textbf{7.} Amplituden $\sqrt{0.9}$, $\sqrt{0.8}$: $V = 2\sqrt{0.72}/1.7 = 0.998$.

\textbf{8.} $(n - 1)e = \qty{5}{\micro m}$: $8.5$ strepen; met natriumlicht
moeiteloos te volgen; in wit licht verschuift de witte streep over
$8.5i = \qty{4.2}{mm}$
en is zij er nog altijd (de dispersie van het glas kost een fractie van een
micrometer).

\textbf{9.} Een witte middenstreep, twee of drie gekleurde strepen aan
elke kant, dan een gelijkmatig wit.

\textbf{10.} Lucht $\to$ olie keert om, olie $\to$ water (lagere index) niet:
$\delta = 2ne + \lambda/2$; helder voor $2ne = (p + \tfrac12)\lambda$.

\textbf{11.} $e = \lambda/4n$: rood \qty{112}{nm}, groen \qty{95}{nm},
violet \qty{72}{nm}.

\textbf{12.} $2ne = \qty{1740}{nm}$: versterkt $1740/(p + \tfrac12)$: $696$
(rood), $497$
(blauwgroen), \qty{387}{nm}; uitgedoofd $1740/p$: $580$ (geel), \qty{435}{nm}:
een purperachtig mengsel van rood en blauwgroen.

\textbf{13.} De dikte verandert langs de weg; elke kleur is helder
waar $2ne$ de juiste waarde heeft, zodat de banden contouren van gelijke
dikte zijn.

\textbf{14.} De orden dalen: de kleuren lopen door de reeks
naar het dunne uiteinde; bij $e \to 0$ laat de enkele tekenomkering
$\delta = \lambda/2$ voor elke kleur over: donker --- zoals bij een
zeepvlies in lucht.

\textbf{15.} Enkele nanometers: $2ne \ll \lambda$ voor alle kleuren, de twee
weerkaatsingen zijn in tegenfase en doven elkaar uit: zwart.

\textbf{16.} $\sin r = \sin60^\circ/1.45$, $r = \qty{37}{\degree}$,
$\cos r = 0.80$: $2ne\cos r
= \qty{1390}{nm}$, helder bij $928$, $557$,
\qty{398}{nm}: het rood van vraag 12 is
groengeel geworden --- naar het blauw toe.

\textbf{17.} Zodra $2ne$ groter wordt dan
$\ell_c \approx \qty{1}{\micro m}$, dat wil zeggen voorbij ongeveer
\qty{0.4}{\micro m} olie, overlappen de orden en vervagen de kleuren tot grijs.

\textbf{18.} De twee stralen die voor welke bronrichting ook door het vlies
worden weerkaatst, ontmoeten elkaar terug op het vlies: de strepen zijn daar
gelokaliseerd en
een brede bron uit vele richtingen voegt slechts gelijke patronen toe.

\textbf{19.} $\ell_c/\lambda$: $34\,000$ met natrium, $1.7 \times 10^7$ met
de laser.

\textbf{20.} $b < \lambda d/4a$: \qty{15}{\micro m} voor de spiegels,
\qty{59}{\micro m} voor
de spleten.

\textbf{21.} Teken twee bronpunten: voor elk komen de twee weerkaatste stralen
in hetzelfde punt van het vlies weer samen, met dezelfde $\delta = 2ne\cos r$
bij bijna loodrechte inval; het patroon van het vlies beweegt niet, dat van Young
wel.

\textbf{22.} $\lambda/\Delta\lambda = 589/0.02 \approx 30\,000$.

\textbf{23.} $550/30 \approx 18$ strepen.

\textbf{24.} Coherent licht dat door de korrels van het scherm wordt
verstrooid, interfereert
met willekeurige fasen in het oog: spikkels.

\textbf{25.} Spiegels: \omterm{prop:b2:two-wave-interference:young}{verdeling van het golffront}, strepen overal
(niet gelokaliseerd), aantal begrensd door $\ell_c/\lambda$, een brede bron wist
ze uit. Vlies: \omterm{prop:b2:two-wave-interference:film}{verdeling van de amplitude}, strepen op het vlies, kleuren
begrensd door $\ell_c \sim \qty{1}{\micro m}$, een brede bron is onschadelijk.
\end{solution}
